\documentclass[a4paper]{article}

\usepackage{graphicx}
\usepackage{amsmath,amssymb}
\usepackage{minted}
\usepackage{multirow}
\usepackage{float}
\usepackage{todonotes}
\usepackage{forest}
\usepackage{xstring}
\usepackage{xcolor}
\usepackage[most]{tcolorbox}
\usepackage{xparse}
\usepackage{natbib}

\usetikzlibrary{graphs,graphdrawing}
\usegdlibrary{layered}

% for external graphviz
\input{/home/donkarlo/Dropbox/repo/nd_language_ai_project/src/nd_language_ai/formal/computer/programming/tex/latex/code_clips/external_graphviz.tex}

% for tags
\input{/home/donkarlo/Dropbox/repo/nd_language_ai_project/src/nd_language_ai/formal/computer/programming/tex/latex/parts/control_sequence/kind/semtag/semtag.tex}

% for links
\input{/home/donkarlo/Dropbox/repo/nd_language_ai_project/src/nd_language_ai/formal/computer/programming/tex/latex/code_clips/seehere.tex}

\title{Ein freund besuchen}
\date{}
\author{Mohammad Rahmani}

\begin{document}
\maketitle
    Liebe Nadin,

    Ich freue mich, dass du seit(immer dativ) einem Jahr(neutral) in Paris lebst.
    
    Die Stadt, in der ich immer gern gelebt hätte
    
    Ich habe vor etwa acht Jahren einmal Paris besucht.
    
    Aber damals hatte ich es sehr eilig.
    
    Es war so, als müsste ich alle Sehenswürdigkeiten besichtigen.
    
    Deshalb stelle ich mir vor, dass du dich inzwischen, da du seit einem Jahr in Paris lebst, gut eingelebt hast und ich ein paar Tage bei dir bleiben kann.

    Überigens,  kannst du bitte mir sagen, wie weit deine Wohnung vom Eiffelturm weg ist?
    
    Ich habe gehört, dass der Eiffelturm nicht mehr der höchste Turm in der Welt ist aber wie beim letzten Mal möchte ich ihn nachts bis zur hochesten erlaubten Stelle besteigen, damit ich die Kalte Luft von Paris auf meinem Gesicht spüren kann.
    
    Eigentlich denke ich schon seit langer Zeit(f), dass ich ein Urlaub(m) weit weg von meinem(Dativ wegen von) regelmäßigen Lebenplatz(m) brauche.
    
    Ich frue mich, dich in meiner Wohnung willkommen(adverb) zu heißen, wann du nach Graz kommen wollen(Modal verb dann vor kommen) wirst (Für Furtur bauen).

    Liebe Grüße
    Mohammad

    
    \bibliographystyle{plainnat}
    \bibliography{/home/donkarlo/Dropbox/repo/nd_shared_research/bibliography/bibliography.bib}
\end{document}